\documentclass[11pt]{article}

% Encoding & fonts (safer defaults)
\usepackage[utf8]{inputenc}
\usepackage[T1]{fontenc}
\usepackage{lmodern}

% Packages for formatting
\usepackage[a4paper,margin=1in]{geometry}
\usepackage{enumitem}
\usepackage{hyperref}
\usepackage{titlesec}
\usepackage{array}
\usepackage{xcolor}
\usepackage{booktabs} % nicer tables
\usepackage{amsmath}  % for math symbols like \rightarrow

% Section formatting
\titleformat{\section}{\large\bfseries}{}{0em}{}
\titleformat{\subsection}{\normalsize\bfseries}{}{0em}{}

% Header info
\newcommand{\university}{University of Vienna}
\newcommand{\coursename}{Applied Economics}
\newcommand{\coursecode}{040178-1 and 040178-2}
\newcommand{\semester}{Winter Term 2026/2027}
\newcommand{\instructor}{Pello Aspuru}
\newcommand{\email}{pello.aspuru@univie.ac.at}

\begin{document}

% Title
\begin{center}
    {\Large \textbf{\coursename: Project Instructions}} \\
    \coursecode \\
    \semester \\
    \vspace{0.3cm}
    {\university} \\
    \vspace{0.3cm}
    Instructor: \instructor \quad Email: \email
\end{center}

\section{Overview}
The project gives you the opportunity to put into practice the skills learnt in class. Working in a team, you will formulate a causal research question in applied economics, collect and explore the data needed to answer it, and estimate a causal effect using one of the methods studied in the course. You will present your work in two short presentations.

\begin{center}
\begin{tabular}{llr}
    \toprule
    \textbf{Item} & \textbf{Date} & \textbf{Grade Weight} \\
    \midrule
    Project Presentation I & December 16th/17th 2026 & 10\% \\
    Project Presentation II & January 20th/21st 2027 & 15\% \\
    \bottomrule
\end{tabular}
\end{center}

\section{Teams}
\begin{itemize}
    \item The project is done in \textbf{teams of three students}.
    \item You have the \textbf{first month of the course} to form your team and register it in the Google Sheet provided on Moodle.
    \item Students who are not registered in a team by then will be assigned to teams \textbf{at random} by the instructor.
    \item If the number of students is not divisible by three, the instructor will create one or two teams of two students.
\end{itemize}

\section{The Four Steps of the Project}

\subsection{Step 1: Research Question}
\begin{itemize}
    \item Come up with a question in applied economics, of the form \textbf{``Does $X$ cause $Y$?''} 
    \begin{itemize}
        \item Do data centers \textbf{cause} an increase in electricity prices? 
        \item Do evangelical curches \textbf{cause} an increase in far-right vote?
    \end{itemize}
    \item Argue why the question is relevant: who cares about the answer, and why?
\end{itemize}

\subsection{Step 2: Data}
\begin{itemize}
    \item Identify several datasets that could be used to answer your question. If neccessary, vibecode a web scraper to create a new dataset or enrich an existing one with new variables.
    \item Understand your data: what is the unit of observation, which period and population are covered, how are the key variables measured?
\end{itemize}

\subsection{Step 3: Graphs}
\begin{itemize}
    \item Provide a preliminary analysis of the data using ggplot2 graphs.
\end{itemize}

\subsection{Step 4: Causal Inference}
\begin{itemize}
    \item Use one of the causal inference methods learnt in class to answer your research question.
\end{itemize}

\section{Presentations}
Each presentation consists of \textbf{2--3 slides (one per team member)} and must last \textbf{no more than 5 minutes in total}.

\subsection{Presentation I (Steps 1--2)}
\begin{enumerate}
    \item Propose your research question and argue about its relevance.
    \item Present the data you plan to use, and how you plan to obtain it.
    \item Argue the reasons why it is the best data to answer your question.
\end{enumerate}

\subsection{Presentation II (Steps 3--4)}
\begin{enumerate}
    \item Pick the most interesting graph you generated with your data, explain why you picked it, and explain the graph to the audience.
    \item Present your empirical specification and the assumptions for a causal interpretation.
    \item Present your parameter estimates and explain them.
\end{enumerate}

\section{Submission}
\begin{itemize}
    \item Upload your presentation to Moodle by \textbf{midnight (23:59) on the day before} your presentation.
    \item \textcolor{red}{Failure to upload on time results in a loss of 5\% of the final course grade for all team members.}
\end{itemize}

\section{Who Presents and Grading}
\begin{itemize}
    \item Every team member must be able to present and answer questions about \textbf{all} slides of the presentation, not only their own.
    \item The instructor will decide on the spot who presents which slide and who answers which question.
    \item Grades are \textbf{individual}, based on each student's performance during the presentation.
\end{itemize}

\section{Academic Integrity}
Students are expected to uphold the highest standards of academic honesty. Plagiarism or cheating will result in disciplinary action in accordance with university policy.

\end{document}
